\clearpage
\fancyhead[L]{\small STERNBEDECKUNG.}
\hypertarget{native-453-sternbedeckung}{}
\begin{center}
{\large STERNBEDECKUNG.}

\rule{0.75\linewidth}{0.4pt}

Astronomische Nachrichten. Band III. Nr. 50. Seite 22. 1824 Februar.

\rule{0.75\linewidth}{0.4pt}
\end{center}

\begin{center}
Bedeckung der Plejaden am 29. August 1820.
\end{center}

Dieselbe Bedeckung ist in Göttingen, Bremen, Altona, Hamburg, Berlin, Wien und Modena beobachtet.

\begin{center}
Göttingen (in einem Briefe von \textsc{Gauss} [an \textsc{Bessel}] gefälligst mitgetheilt).

\begingroup\small\setlength{\tabcolsep}{4pt}
\begin{tabular}{c l c c c | c c l}
$p$ & Austritt & $9^{\mathrm u}$ & $46^{\mathrm m}$ & $42^{\mathrm s}49$ & $10^{\mathrm u}$ & $5^{\mathrm m}$ & $26^{\mathrm s}78 - 1.582x$ \\
$\eta$ & --- & $9$ & $51$ & $37.37$ & $10$ & $8$ & $22.89 - 1.265x$ \\
$f$ & Eintritt & $9$ & $40$ & $4.77$ & $10$ & $46$ & $56.01 + 0.219x$ \\
$h$ & --- & $9$ & $44$ & $11.60$ & $10$ & $49$ & $18.09 + 0.795x$ \\
\end{tabular}
\endgroup
\end{center}

[wo $x$ die Verbesserung der Breitendifferenz des Mondes und Sterns bezeichnet.]

\clearpage
\fancyhead[L]{\small GAUSS AN BODE.}
\SourceLeaf{454}

\begin{center}
{\large GAUSS AN BODE.}
\end{center}

\begin{center}
\rule{6cm}{0.4pt}\\
Astronomisches Jahrbuch für 1827. Seite 128..130. Berlin 1824.\\
\rule{6cm}{0.4pt}
\end{center}

\begin{flushright}
Göttingen den 21. April 1824.
\end{flushright}

Bei allen diesen Beobachtungen war der Comet so lichtschwach, dass er gar keine
Fadenbeleuchtung vertrug, Comet und Fadennetz konnten daher nur bei
abwechselndem Schliessen und Oeffnen der Beleuchtung sichtbar gemacht werden.
Unter diesen Umständen waren die Beobachtungen, besonders die Rectascensionen,
nur geringer Genauigkeit fähig. Die beigesetzten Zeiten beziehen sich auf das
Mittel der beobachteten oder vielmehr geschätzten Antritte, und sind also an
solchen Tagen, wo nicht alle Faden beobachtet wurden, von der wahren
Culminationszeit etwas verschieden.

\begin{center}
\begingroup\small\setlength{\tabcolsep}{4pt}
\begin{tabular}{lr|c|c|c}
\multicolumn{2}{c|}{1824} & \shortstack{Mittlere Zeit\\in Göttingen} & \shortstack{Gerade\\Aufsteigung} & Abweichung \\
\hline
Febr. & 31 & $13^{\mathrm u}\,28^{\mathrm m}\,20^{\mathrm s}$ & $152^\circ\,16'\,29''$ & $72^\circ\,38'\,35''$ N. \\
Jan. & 3 & $11\quad55\quad7$ & $131\quad23\quad10$ & $69\quad14\quad9$ \\
 & 7 & $10\quad38\quad3$ & $116\quad12\quad50$ & $63\quad14\quad56$ \\
 & 8 & $10\quad24\quad5$ & $113\quad50\quad49$ & $61\quad46\quad15$ \\
 & 19 & $8\quad52\quad23$ & $101\quad49\quad31$ & $48\quad48\quad51$ \\
 & 20 & $8\quad48$ & ---\quad---\quad--- & $47\quad55$ \quad : \\
 & 27 & $8\quad11\quad46$ & $99\quad26\quad51$ & $42\quad48\quad49$ \\
 & 28 & $8\quad6\quad47$ & $99\quad18\quad8$ & $42\quad12\quad46$ \\
März & 2 & $7\quad53\quad59$ & $99\quad0\quad39$ & $40\quad32\quad3$ \\
\end{tabular}
\endgroup
\end{center}

Einige frühere ausser dem Meridian, mit Kreismikrometer oder Heliometer von mir
gemachte Beobachtungen, habe ich noch nicht reducirt, da ich die in der
Histoire Céleste nicht vorkommenden Vergleichungs-Sterne erst selbst bestimmen
muss.

Hr. Dr. \textsc{Schmidt}, ein junger Mann von ausgezeichneten mathematischen
Kenntnissen, hat nach diesen und einigen andern Beobachtungen folgende
parabolische Elemente berechnet:

\begin{center}
\begingroup\small
\begin{tabular}{p{.51\linewidth}p{.41\linewidth}}
Länge des Knoten \mbox{.\,.\,.\,.\,.\,.\,.\,.\,.\,.\,.\,.} & $303^\circ\,3'\,51''3$ \\
Neigung der Bahn \mbox{.\,.\,.\,.\,.\,.\,.\,.\,.\,.\,.\,.} & $76\quad12\quad14.0$ \\
Länge der Sonnennähe \mbox{.\,.\,.\,.\,.\,.\,.\,.\,.\,.\,.} & $274\quad34\quad14.4$ \\
Durchgang durch dieselbe \mbox{.\,.\,.\,.} & 1823 December 9.47459 Göttinger Zeit. \\
Logarithm des kleinsten Abstandes. \mbox{.\,.\,.\,.\,.\,.} & 9.3553041 \\
\multicolumn{2}{l}{Bewegung rückläufig.} \\
\end{tabular}
\endgroup
\end{center}

Die Vergleichung dieser Elemente mit den bisher bekannt gewordenen
Beobachtungen gab folgende Unterschiede:

\Needspace{12\baselineskip}
\SourceLeaf{455}


\begin{center}
\begin{tabular}{lr|r|r|l}
 & & \shortstack{in Gerader\\Aufsteig.} & \shortstack{in\\Abweichung} & Beobachter \\
\hline
Jan. & 1 & 0 & 0 & Bouvard \\
 & 3 & $+12.5$ & $-22.6$ & Nicolai \\
 & 4 & $+13.6$ & $+1.6$ & Nicolai \\
 & 5 & $+34.5$ & $-6.5$ & Nicolai \\
 & 11 & $+0.2$ & $+9.5$ & Olbers \\
 & 14 & $+27.5$ & $+0.3$ & Olbers \\
 & 24 & $+48.6$ & $+3.4$ & Schumacher \\
 & 25 & $+85.0$ & $-3.3$ & Schumacher \\
 & 26 & $+44.0$ & $+11.0$ & Soldner \\
 & 27 & $+33.0$ & $+10.4$ & Soldner \\
 & 30 & $-26.5$ & $+20.7$ & Schumacher \\
 & 31 & $-36.7$ & $+25.5$ & Soldner \\
 & 31 & $+24.9$ & $+14.5$ & Gauss \\
Febr. & 2 & $+963.4$ & $-5.7$ & Schumacher \\
 & 3 & $-28.9$ & $+12.0$ & Schumacher \\
 & 3 & $-40.6$ & $-8.6$ & Gauss \\
\end{tabular}
\end{center}

\begin{center}
\begin{tabular}{lr|r|r|l}
 & & \shortstack{in Gerader\\Aufsteig.} & \shortstack{in\\Abweichung} & Beobachter \\
\hline
Febr. & 7 & $-4.5$ & $+18.5$ & Gauss \\
 & 8 & $-40.6$ & $-2.0$ & Gauss \\
 & 11 & $-18.4$ & $-22.0$ & Schumacher \\
 & 13 & $-10.4$ & $-1.2$ & Schumacher \\
 & 18 & $-2.6$ & $-66.1$ & Olbers \\
 & 19 & $-1.0$ & $+1.0$ & Gauss \\
 & 19 & $+26.0$ & $-26.7$ & Olbers \\
 & 21 & $+36.2$ & $-38.8$ & Olbers \\
 & 27 & $+27.8$ & $-0.2$ & Gauss \\
 & 27 & $+21.6$ & $-34.5$ & Olbers \\
 & 28 & $+21.4$ & $-22.7$ & Gauss \\
 & 28 & $+0.4$ & $-11.4$ & Olbers \\
März & 2 & $-41.3$ & $+9.8$ & Gauss \\
 & 2 & $+40.1$ & $-34.9$ & Olbers \\
 & 5 & $+7.3$ & $-83.3$ & Olbers \\
 & 19 & $+56.3$ & $+36.4$ & Olbers \\
\end{tabular}
\end{center}


Bei der Rectascension der Altonaer Beobachtung vom 2. Februar scheint ein
Fehler von Einer Zeit-Minute vorgefallen zu sein.
